\documentclass[8pt,a4paper]{extarticle}
% ============================================================
% Gemeinsamer Preamble fuer alle 5 Open-Book-Spickzettel.
% Wird aus jedem klausur.tex via % ============================================================
% Gemeinsamer Preamble fuer alle 5 Open-Book-Spickzettel.
% Wird aus jedem klausur.tex via % ============================================================
% Gemeinsamer Preamble fuer alle 5 Open-Book-Spickzettel.
% Wird aus jedem klausur.tex via \input{../preamble.tex} geladen.
% Layout: A4 hochkant, 3 Spalten, 8pt, sehr kompakt.
% ============================================================

\usepackage[a4paper,margin=8mm,top=7mm,bottom=7mm,includefoot]{geometry}
\usepackage[utf8]{inputenc}
\usepackage[T1]{fontenc}
\usepackage[ngerman]{babel}
\usepackage{lmodern}
\usepackage{microtype}

\usepackage{amsmath,amssymb,amsfonts,mathtools}
\usepackage{multicol}
\usepackage{enumitem}
\usepackage{array,tabularx,booktabs,makecell}
\usepackage[dvipsnames,table]{xcolor}
\usepackage[explicit]{titlesec}
\usepackage{etoolbox}
\usepackage{fancyhdr}
\usepackage{lastpage}

% ---------- Farben ----------
\definecolor{secblue}{HTML}{1f4e79}
\definecolor{subblue}{HTML}{2e75b6}
\definecolor{boxbg}{HTML}{f2f2f2}
\definecolor{accent}{HTML}{c00000}
\definecolor{rulegray}{HTML}{888888}

% ---------- Spalten ----------
\setlength{\columnsep}{3mm}
\setlength{\columnseprule}{0.3pt}
\def\columnseprulecolor{\color{rulegray}}

% ---------- Listen kompakt ----------
\setlist{nosep,leftmargin=1.1em,labelsep=0.35em,topsep=0.2ex,partopsep=0pt}
\setlist[itemize]{label=\textbullet}

% ---------- Abstaende ----------
\setlength{\parindent}{0pt}
\setlength{\parskip}{0.15ex}
\setlength{\abovedisplayskip}{0.3ex}
\setlength{\belowdisplayskip}{0.3ex}
\setlength{\abovedisplayshortskip}{0.2ex}
\setlength{\belowdisplayshortskip}{0.2ex}
\setlength{\jot}{1pt}

% ---------- Tabellen kompakt ----------
\renewcommand{\arraystretch}{1.05}
\setlength{\tabcolsep}{3pt}
\setlength{\arrayrulewidth}{0.3pt}
\arrayrulecolor{rulegray}

% ---------- Section-Stil (kein Vollblock; dick + farbige Linie darunter) ----------
\titleformat{\section}
  {\large\bfseries\color{secblue}}
  {}{0pt}{#1}
  [\vspace{-0.4ex}{\color{secblue}\hrule height 0.8pt}\vspace{0.1ex}]
\titlespacing*{\section}{0pt}{1.4ex plus 0.4ex}{0.5ex}

\titleformat{\subsection}
  {\small\bfseries\color{subblue}}
  {}{0pt}{#1}
\titlespacing*{\subsection}{0pt}{0.9ex plus 0.2ex}{0.2ex}

\titleformat{\subsubsection}
  {\small\bfseries}
  {}{0pt}{#1}
\titlespacing*{\subsubsection}{0pt}{0.5ex}{0.1ex}

% ---------- Mini-Helfer ----------

% \fact{...} Hervorgehobene Regel (gerahmt, hellgrau)
\newcommand{\fact}[1]{%
  \begingroup\setlength{\fboxsep}{2pt}%
  \colorbox{boxbg}{\parbox{\dimexpr\linewidth-2\fboxsep}{\small #1\strut}}%
  \endgroup\par\vspace{0.2ex}}

% \warn{...} Stolperfalle (roter Strich links)
\newcommand{\warn}[1]{%
  \begingroup\setlength{\fboxsep}{2pt}%
  \noindent{\color{accent}\rule{2pt}{\baselineskip}}\hspace{2pt}%
  \begin{minipage}[t]{\dimexpr\linewidth-6pt}\small #1\end{minipage}%
  \endgroup\par\vspace{0.2ex}}

% Math-Shortcuts
\newcommand{\R}{\mathbb{R}}
\newcommand{\N}{\mathbb{N}}
\newcommand{\Z}{\mathbb{Z}}
\newcommand{\Q}{\mathbb{Q}}
\newcommand{\C}{\mathbb{C}}
\DeclareMathOperator{\rg}{rg}
\DeclareMathOperator{\spur}{tr}
\DeclareMathOperator{\im}{im}
\DeclareMathOperator{\diag}{diag}
\DeclareMathOperator{\sign}{sgn}
\DeclareMathOperator*{\argmin}{arg\,min}
\DeclareMathOperator*{\argmax}{arg\,max}
\newcommand{\trans}{^{\!\top}}
\newcommand{\norm}[1]{\left\lVert #1 \right\rVert}
\newcommand{\abs}[1]{\left\lvert #1 \right\rvert}
\newcommand{\inner}[2]{\langle #1,\,#2\rangle}
\newcommand{\dx}{\,\mathrm{d}x}
\newcommand{\dy}{\,\mathrm{d}y}
\newcommand{\dt}{\,\mathrm{d}t}

% Zeigt eine Display-Formel zentriert; skaliert nur runter wenn zu breit.
\newlength{\fitw}
\newcommand{\fitline}[1]{%
  \par\nointerlineskip
  \settowidth{\fitw}{$\displaystyle #1$}%
  \ifdim\fitw>\linewidth
    \noindent\resizebox{\linewidth}{!}{$\displaystyle #1$}\par
  \else
    \[\displaystyle #1\]%
  \fi}

% Tabular fuer fixe Spaltenbreiten (Spickzettel-typisch)
\newenvironment{tightt}[1]
  {\par\noindent\begin{tabular}{#1}}
  {\end{tabular}\par}
% Fuer X-Spalten direkt \begin{tabularx}{\linewidth}{...} verwenden
% (kann nicht in \newenvironment gewrappt werden -- TX@get@body Konflikt)

% Globale Schutzschalter gegen Overflow
\emergencystretch=2em
\tolerance=2000
\hbadness=10000
\hfuzz=2pt
\vfuzz=2pt

% ---------- Kopf-/Fusszeile ----------
\pagestyle{fancy}
\fancyhf{}
\renewcommand{\headrulewidth}{0pt}
\renewcommand{\footrulewidth}{0pt}
\fancyfoot[L]{\scriptsize\color{rulegray}\textit{\sheettitle}}
\fancyfoot[R]{\scriptsize\color{rulegray}Seite \thepage\,/\,\pageref{LastPage}}
\setlength{\footskip}{8pt}

% Titel-Header oben auf Seite 1
\newcommand{\sheetheader}[2]{%
  {\Large\bfseries\color{secblue} #1}\hfill
  {\small\color{rulegray} #2}\par
  \vspace{0.4ex}{\color{secblue}\hrule height 0.7pt}\vspace{1ex}}

% Default-Titel (wird ueberschrieben)
\newcommand{\sheettitle}{Open-Book Spickzettel}
 geladen.
% Layout: A4 hochkant, 3 Spalten, 8pt, sehr kompakt.
% ============================================================

\usepackage[a4paper,margin=8mm,top=7mm,bottom=7mm,includefoot]{geometry}
\usepackage[utf8]{inputenc}
\usepackage[T1]{fontenc}
\usepackage[ngerman]{babel}
\usepackage{lmodern}
\usepackage{microtype}

\usepackage{amsmath,amssymb,amsfonts,mathtools}
\usepackage{multicol}
\usepackage{enumitem}
\usepackage{array,tabularx,booktabs,makecell}
\usepackage[dvipsnames,table]{xcolor}
\usepackage[explicit]{titlesec}
\usepackage{etoolbox}
\usepackage{fancyhdr}
\usepackage{lastpage}

% ---------- Farben ----------
\definecolor{secblue}{HTML}{1f4e79}
\definecolor{subblue}{HTML}{2e75b6}
\definecolor{boxbg}{HTML}{f2f2f2}
\definecolor{accent}{HTML}{c00000}
\definecolor{rulegray}{HTML}{888888}

% ---------- Spalten ----------
\setlength{\columnsep}{3mm}
\setlength{\columnseprule}{0.3pt}
\def\columnseprulecolor{\color{rulegray}}

% ---------- Listen kompakt ----------
\setlist{nosep,leftmargin=1.1em,labelsep=0.35em,topsep=0.2ex,partopsep=0pt}
\setlist[itemize]{label=\textbullet}

% ---------- Abstaende ----------
\setlength{\parindent}{0pt}
\setlength{\parskip}{0.15ex}
\setlength{\abovedisplayskip}{0.3ex}
\setlength{\belowdisplayskip}{0.3ex}
\setlength{\abovedisplayshortskip}{0.2ex}
\setlength{\belowdisplayshortskip}{0.2ex}
\setlength{\jot}{1pt}

% ---------- Tabellen kompakt ----------
\renewcommand{\arraystretch}{1.05}
\setlength{\tabcolsep}{3pt}
\setlength{\arrayrulewidth}{0.3pt}
\arrayrulecolor{rulegray}

% ---------- Section-Stil (kein Vollblock; dick + farbige Linie darunter) ----------
\titleformat{\section}
  {\large\bfseries\color{secblue}}
  {}{0pt}{#1}
  [\vspace{-0.4ex}{\color{secblue}\hrule height 0.8pt}\vspace{0.1ex}]
\titlespacing*{\section}{0pt}{1.4ex plus 0.4ex}{0.5ex}

\titleformat{\subsection}
  {\small\bfseries\color{subblue}}
  {}{0pt}{#1}
\titlespacing*{\subsection}{0pt}{0.9ex plus 0.2ex}{0.2ex}

\titleformat{\subsubsection}
  {\small\bfseries}
  {}{0pt}{#1}
\titlespacing*{\subsubsection}{0pt}{0.5ex}{0.1ex}

% ---------- Mini-Helfer ----------

% \fact{...} Hervorgehobene Regel (gerahmt, hellgrau)
\newcommand{\fact}[1]{%
  \begingroup\setlength{\fboxsep}{2pt}%
  \colorbox{boxbg}{\parbox{\dimexpr\linewidth-2\fboxsep}{\small #1\strut}}%
  \endgroup\par\vspace{0.2ex}}

% \warn{...} Stolperfalle (roter Strich links)
\newcommand{\warn}[1]{%
  \begingroup\setlength{\fboxsep}{2pt}%
  \noindent{\color{accent}\rule{2pt}{\baselineskip}}\hspace{2pt}%
  \begin{minipage}[t]{\dimexpr\linewidth-6pt}\small #1\end{minipage}%
  \endgroup\par\vspace{0.2ex}}

% Math-Shortcuts
\newcommand{\R}{\mathbb{R}}
\newcommand{\N}{\mathbb{N}}
\newcommand{\Z}{\mathbb{Z}}
\newcommand{\Q}{\mathbb{Q}}
\newcommand{\C}{\mathbb{C}}
\DeclareMathOperator{\rg}{rg}
\DeclareMathOperator{\spur}{tr}
\DeclareMathOperator{\im}{im}
\DeclareMathOperator{\diag}{diag}
\DeclareMathOperator{\sign}{sgn}
\DeclareMathOperator*{\argmin}{arg\,min}
\DeclareMathOperator*{\argmax}{arg\,max}
\newcommand{\trans}{^{\!\top}}
\newcommand{\norm}[1]{\left\lVert #1 \right\rVert}
\newcommand{\abs}[1]{\left\lvert #1 \right\rvert}
\newcommand{\inner}[2]{\langle #1,\,#2\rangle}
\newcommand{\dx}{\,\mathrm{d}x}
\newcommand{\dy}{\,\mathrm{d}y}
\newcommand{\dt}{\,\mathrm{d}t}

% Zeigt eine Display-Formel zentriert; skaliert nur runter wenn zu breit.
\newlength{\fitw}
\newcommand{\fitline}[1]{%
  \par\nointerlineskip
  \settowidth{\fitw}{$\displaystyle #1$}%
  \ifdim\fitw>\linewidth
    \noindent\resizebox{\linewidth}{!}{$\displaystyle #1$}\par
  \else
    \[\displaystyle #1\]%
  \fi}

% Tabular fuer fixe Spaltenbreiten (Spickzettel-typisch)
\newenvironment{tightt}[1]
  {\par\noindent\begin{tabular}{#1}}
  {\end{tabular}\par}
% Fuer X-Spalten direkt \begin{tabularx}{\linewidth}{...} verwenden
% (kann nicht in \newenvironment gewrappt werden -- TX@get@body Konflikt)

% Globale Schutzschalter gegen Overflow
\emergencystretch=2em
\tolerance=2000
\hbadness=10000
\hfuzz=2pt
\vfuzz=2pt

% ---------- Kopf-/Fusszeile ----------
\pagestyle{fancy}
\fancyhf{}
\renewcommand{\headrulewidth}{0pt}
\renewcommand{\footrulewidth}{0pt}
\fancyfoot[L]{\scriptsize\color{rulegray}\textit{\sheettitle}}
\fancyfoot[R]{\scriptsize\color{rulegray}Seite \thepage\,/\,\pageref{LastPage}}
\setlength{\footskip}{8pt}

% Titel-Header oben auf Seite 1
\newcommand{\sheetheader}[2]{%
  {\Large\bfseries\color{secblue} #1}\hfill
  {\small\color{rulegray} #2}\par
  \vspace{0.4ex}{\color{secblue}\hrule height 0.7pt}\vspace{1ex}}

% Default-Titel (wird ueberschrieben)
\newcommand{\sheettitle}{Open-Book Spickzettel}
 geladen.
% Layout: A4 hochkant, 3 Spalten, 8pt, sehr kompakt.
% ============================================================

\usepackage[a4paper,margin=8mm,top=7mm,bottom=7mm,includefoot]{geometry}
\usepackage[utf8]{inputenc}
\usepackage[T1]{fontenc}
\usepackage[ngerman]{babel}
\usepackage{lmodern}
\usepackage{microtype}

\usepackage{amsmath,amssymb,amsfonts,mathtools}
\usepackage{multicol}
\usepackage{enumitem}
\usepackage{array,tabularx,booktabs,makecell}
\usepackage[dvipsnames,table]{xcolor}
\usepackage[explicit]{titlesec}
\usepackage{etoolbox}
\usepackage{fancyhdr}
\usepackage{lastpage}

% ---------- Farben ----------
\definecolor{secblue}{HTML}{1f4e79}
\definecolor{subblue}{HTML}{2e75b6}
\definecolor{boxbg}{HTML}{f2f2f2}
\definecolor{accent}{HTML}{c00000}
\definecolor{rulegray}{HTML}{888888}

% ---------- Spalten ----------
\setlength{\columnsep}{3mm}
\setlength{\columnseprule}{0.3pt}
\def\columnseprulecolor{\color{rulegray}}

% ---------- Listen kompakt ----------
\setlist{nosep,leftmargin=1.1em,labelsep=0.35em,topsep=0.2ex,partopsep=0pt}
\setlist[itemize]{label=\textbullet}

% ---------- Abstaende ----------
\setlength{\parindent}{0pt}
\setlength{\parskip}{0.15ex}
\setlength{\abovedisplayskip}{0.3ex}
\setlength{\belowdisplayskip}{0.3ex}
\setlength{\abovedisplayshortskip}{0.2ex}
\setlength{\belowdisplayshortskip}{0.2ex}
\setlength{\jot}{1pt}

% ---------- Tabellen kompakt ----------
\renewcommand{\arraystretch}{1.05}
\setlength{\tabcolsep}{3pt}
\setlength{\arrayrulewidth}{0.3pt}
\arrayrulecolor{rulegray}

% ---------- Section-Stil (kein Vollblock; dick + farbige Linie darunter) ----------
\titleformat{\section}
  {\large\bfseries\color{secblue}}
  {}{0pt}{#1}
  [\vspace{-0.4ex}{\color{secblue}\hrule height 0.8pt}\vspace{0.1ex}]
\titlespacing*{\section}{0pt}{1.4ex plus 0.4ex}{0.5ex}

\titleformat{\subsection}
  {\small\bfseries\color{subblue}}
  {}{0pt}{#1}
\titlespacing*{\subsection}{0pt}{0.9ex plus 0.2ex}{0.2ex}

\titleformat{\subsubsection}
  {\small\bfseries}
  {}{0pt}{#1}
\titlespacing*{\subsubsection}{0pt}{0.5ex}{0.1ex}

% ---------- Mini-Helfer ----------

% \fact{...} Hervorgehobene Regel (gerahmt, hellgrau)
\newcommand{\fact}[1]{%
  \begingroup\setlength{\fboxsep}{2pt}%
  \colorbox{boxbg}{\parbox{\dimexpr\linewidth-2\fboxsep}{\small #1\strut}}%
  \endgroup\par\vspace{0.2ex}}

% \warn{...} Stolperfalle (roter Strich links)
\newcommand{\warn}[1]{%
  \begingroup\setlength{\fboxsep}{2pt}%
  \noindent{\color{accent}\rule{2pt}{\baselineskip}}\hspace{2pt}%
  \begin{minipage}[t]{\dimexpr\linewidth-6pt}\small #1\end{minipage}%
  \endgroup\par\vspace{0.2ex}}

% Math-Shortcuts
\newcommand{\R}{\mathbb{R}}
\newcommand{\N}{\mathbb{N}}
\newcommand{\Z}{\mathbb{Z}}
\newcommand{\Q}{\mathbb{Q}}
\newcommand{\C}{\mathbb{C}}
\DeclareMathOperator{\rg}{rg}
\DeclareMathOperator{\spur}{tr}
\DeclareMathOperator{\im}{im}
\DeclareMathOperator{\diag}{diag}
\DeclareMathOperator{\sign}{sgn}
\DeclareMathOperator*{\argmin}{arg\,min}
\DeclareMathOperator*{\argmax}{arg\,max}
\newcommand{\trans}{^{\!\top}}
\newcommand{\norm}[1]{\left\lVert #1 \right\rVert}
\newcommand{\abs}[1]{\left\lvert #1 \right\rvert}
\newcommand{\inner}[2]{\langle #1,\,#2\rangle}
\newcommand{\dx}{\,\mathrm{d}x}
\newcommand{\dy}{\,\mathrm{d}y}
\newcommand{\dt}{\,\mathrm{d}t}

% Zeigt eine Display-Formel zentriert; skaliert nur runter wenn zu breit.
\newlength{\fitw}
\newcommand{\fitline}[1]{%
  \par\nointerlineskip
  \settowidth{\fitw}{$\displaystyle #1$}%
  \ifdim\fitw>\linewidth
    \noindent\resizebox{\linewidth}{!}{$\displaystyle #1$}\par
  \else
    \[\displaystyle #1\]%
  \fi}

% Tabular fuer fixe Spaltenbreiten (Spickzettel-typisch)
\newenvironment{tightt}[1]
  {\par\noindent\begin{tabular}{#1}}
  {\end{tabular}\par}
% Fuer X-Spalten direkt \begin{tabularx}{\linewidth}{...} verwenden
% (kann nicht in \newenvironment gewrappt werden -- TX@get@body Konflikt)

% Globale Schutzschalter gegen Overflow
\emergencystretch=2em
\tolerance=2000
\hbadness=10000
\hfuzz=2pt
\vfuzz=2pt

% ---------- Kopf-/Fusszeile ----------
\pagestyle{fancy}
\fancyhf{}
\renewcommand{\headrulewidth}{0pt}
\renewcommand{\footrulewidth}{0pt}
\fancyfoot[L]{\scriptsize\color{rulegray}\textit{\sheettitle}}
\fancyfoot[R]{\scriptsize\color{rulegray}Seite \thepage\,/\,\pageref{LastPage}}
\setlength{\footskip}{8pt}

% Titel-Header oben auf Seite 1
\newcommand{\sheetheader}[2]{%
  {\Large\bfseries\color{secblue} #1}\hfill
  {\small\color{rulegray} #2}\par
  \vspace{0.4ex}{\color{secblue}\hrule height 0.7pt}\vspace{1ex}}

% Default-Titel (wird ueberschrieben)
\newcommand{\sheettitle}{Open-Book Spickzettel}

\renewcommand{\sheettitle}{Lineare Algebra}

\begin{document}
\sheetheader{Lineare Algebra \& Matrix-Zerlegungen}{Open-Book Formelsammlung -- DHBW WDS125}

\begin{multicols*}{3}

% =================================================
\section{Grundlagen}
% =================================================
\subsection{Wichtige Identit\"aten}
\begin{itemize}
\item $\rg(A){=}\dim\im(A){=}\dim$ Zeilen-/Spaltenraum
\item Dimensionssatz: $\dim\ker(A)+\rg(A)=n$
\item $A^{-1}$ existiert $\Leftrightarrow \det A\!\ne\!0\Leftrightarrow \rg(A){=}n\Leftrightarrow$ Spalten lin.\ unabh.
\item $\det(AB){=}\det A\cdot\det B$,\ $\det(A\trans){=}\det A$,\ $\det(A^{-1}){=}1/\det A$
\item $\spur(A){=}\sum_i A_{ii}{=}\sum_i \lambda_i$
\item $\det A=\prod_i\lambda_i$
\item Char.\ Pol.: $\chi_A(\lambda){=}\det(A-\lambda I)$, Grad $n$
\end{itemize}

\subsection{Zerlegungs\"ubersicht}
\begin{tabularx}{\linewidth}{@{}l@{\ }X@{}}
\textbf{LU} & quadr., $\det\!\ne\!0$ -- $Ax{=}b$ \\
\textbf{Cholesky} $LL\trans$ & symm.\ \& PD -- halber LU \\
\textbf{QR} & $m\!\times\!n$, $m\!\ge\!n$ -- Least Sq.\ \\
\textbf{EVD} $PDP^{-1}$ & diagonalisierbar -- $A^k$, ODE \\
\textbf{Jordan} & jede quadr. \\
\textbf{SVD} $U\Sigma V\trans$ & jede $m\!\times\!n$ -- Pseudoinv., PCA, $\kappa$ \\
\end{tabularx}

% =================================================
\section{LU-Zerlegung}
% =================================================
$A=LU$ ($L$ unt.\ $\triangle$, Diag $1$; $U$ ob.\ $\triangle$). Mit Pivot: $PA=LU$. Aufwand $\tfrac23 n^3$.\\
\textbf{Algorithmus} (ohne Pivot):
\begin{enumerate}
\item F\"ur Spalte $j$: f\"ur jede Zeile $i>j$: $\ell_{ij}=a_{ij}/a_{jj}$.
\item $Z_i\leftarrow Z_i-\ell_{ij}Z_j$.
\item \"Ubrig: $U$; Multiplikatoren bilden $L$.
\end{enumerate}
\textbf{$Ax=b$ via LU}: (1) $Ly=b$ vorw\"arts; (2) $Ux=y$ r\"uckw\"arts. Beides $O(n^2)$.
\warn{Spaltenpivotsuche, wenn $a_{jj}$ klein/null.}

% =================================================
\section{Cholesky}
% =================================================
$A=LL\trans$ f\"ur symm.\ \& PD; Aufwand $\tfrac13 n^3$.
\fitline{\ell_{jj}=\sqrt{a_{jj}-\textstyle\sum_{k=1}^{j-1}\ell_{jk}^2}}
\fitline{\ell_{ij}=\tfrac{1}{\ell_{jj}}\Bigl(a_{ij}-\textstyle\sum_{k=1}^{j-1}\ell_{ik}\ell_{jk}\Bigr),\ i>j}
\warn{$\sqrt{\text{negativ}}\Rightarrow A$ nicht PD.}

% =================================================
\section{QR-Zerlegung}
% =================================================
$A=QR$, $Q\trans Q=I$, $R$ ob.\ $\triangle$.\\
\textbf{Klassisches Gram-Schmidt}:
\begin{enumerate}
\item $u_1=a_1$, $q_1=u_1/\|u_1\|$.
\item $u_k=a_k-\sum_{j<k}\inner{a_k}{q_j}q_j$, $q_k=u_k/\|u_k\|$.
\item $R_{jk}=\inner{a_k}{q_j}$ f\"ur $j<k$, $R_{kk}=\|u_k\|$.
\end{enumerate}
\warn{Klass.\ GS instabil $\to$ MGS oder Householder.}
\fact{\textbf{Least Squares}: $Rx=Q\trans b$ statt $A\trans A x=A\trans b$ ($\kappa(A\trans A){=}\kappa(A)^2$).}

% =================================================
\section{Eigenwerte / EVD}
% =================================================
\subsection{Definitionen}
$Av=\lambda v$, $v\ne 0$. Char.\ Pol.\ $\chi_A(\lambda){=}0$, Eigenraum $E_\lambda{=}\ker(A-\lambda I)$.\\
Alg.\ Vielfachheit (in $\chi_A$); geom.\ Vielfachheit $=\dim E_\lambda$;\ $1\!\le\!\text{geom}\!\le\!\text{alg}$.

\subsection{Rezept}
\begin{enumerate}
\item $\det(A-\lambda I){=}0 \Rightarrow \lambda_i$.
\item $(A-\lambda_i I)v=0 \Rightarrow$ Eigenvektoren.
\item Geom.\ Vielf.\ $=$ \#lin.\ unabh.\ L\"osungen.
\end{enumerate}

\subsection{Diagonalisierung}
$A$ diag.\ $\Leftrightarrow$ $\sum$ geom.\ $=n$ (geom\,=\,alg f\"ur jeden EW).\\
$A=PDP^{-1}$, $D=\diag(\lambda_i)$, Spalten von $P$ = EV.\\
$n$ versch.\ EW $\Rightarrow$ diag.\ \\
Symm.\ reell $\Rightarrow$ orthogonal diag.: $A=Q\Lambda Q\trans$.\\
Normal ($AA\trans{=}A\trans A$) $\Rightarrow$ unit\"ar diag.\\
\fact{Anwendung: $A^k=PD^kP^{-1}$.}

\subsection{Jordan-Normalform}
$A=PJP^{-1}$, $J=\diag(J_{k_i}(\lambda_i))$ mit
\[ J_3(\lambda)=\begin{pmatrix}\lambda&1&0\\0&\lambda&1\\0&0&\lambda\end{pmatrix} \]
Blockgr\"ossen aus $\dim\ker(A-\lambda I)^k$-Folge.

% =================================================
\section{SVD}
% =================================================
$A=U\Sigma V\trans$, $A\in\R^{m\times n}$.
\begin{itemize}
\item $U$ ($m{\times}m$) orth.\ -- Links-Sing.-V.
\item $V$ ($n{\times}n$) orth.\ -- Rechts-Sing.-V.
\item $\Sigma$ ``diag.'' mit $\sigma_1\!\ge\!\dots\!\ge\!\sigma_r>0$, $r{=}\rg(A)$.
\end{itemize}
\fact{$A\trans A{=}V\Sigma\trans\Sigma V\trans$,\ $AA\trans{=}U\Sigma\Sigma\trans U\trans$. $\sigma_i^2$=EW von $A\trans A$.}

\subsection{Rezept}
\begin{enumerate}
\item $A\trans A$ bilden.
\item EW $\lambda_i\Rightarrow\sigma_i=\sqrt{\lambda_i}$, absteigend.
\item Orthon.\ EV $v_i$ $\Rightarrow$ Spalten von $V$.
\item $u_i=Av_i/\sigma_i$ f\"ur $\sigma_i>0$.
\item Rest von $U$ aus $\ker(A\trans)$ erg\"anzen (orthonorm.).
\end{enumerate}

\subsection{Anwendungen}
\textbf{Eckart-Young} (beste Rang-$k$-Approx.): $A_k=\sum_{i=1}^k\sigma_i u_i v_i\trans$
\fitline{\|A-A_k\|_2=\sigma_{k+1},\ \|A-A_k\|_F=\sqrt{\textstyle\sum_{i>k}\sigma_i^2}}
\textbf{Pseudoinverse}: $A^+=V\Sigma^+U\trans$, $\Sigma^+_{ii}=1/\sigma_i$ wenn $\sigma_i\!>\!0$, sonst $0$; $\hat x=A^+ b$.\\
\textbf{Konditionszahl}: $\kappa_2(A)=\sigma_1/\sigma_n$.\\
\textbf{PCA}: zentrierte Datenmatrix; $V$ = Hauptachsen, $\sigma_i^2/(m{-}1)$ = Varianz.

% =================================================
\section{Positiv-Definitheit}
% =================================================
F\"ur symm.\ $A$ \"aquivalent:
\begin{enumerate}
\item $x\trans Ax>0\ \forall x\ne 0$
\item alle EW $>0$
\item alle f\"uhrenden Hauptminoren $>0$ (\textbf{Sylvester})
\item Cholesky mit pos.\ Diag.\ existiert
\end{enumerate}
PSD: $\ge 0$ in allen vier.\\
2$\times$2 PD: $a>0$ \emph{und} $ac-b^2>0$.

% =================================================
\section{Geometrie}
% =================================================
\subsection{Skalarprod.\ \& Norm}
$\inner{u}{v}=u\trans v=\sum u_iv_i$;\ $\|v\|=\sqrt{\inner{v}{v}}$.\\
$\cos\theta=\inner{u}{v}/(\|u\|\|v\|)$;\ $u\perp v\Leftrightarrow\inner{u}{v}=0$.\\
\textbf{Cauchy-Schwarz}: $|\inner{u}{v}|\le\|u\|\|v\|$.

\subsection{Kreuzprodukt ($\R^3$)}
\fitline{u\times v=\bigl(u_2v_3{-}u_3v_2,\,u_3v_1{-}u_1v_3,\,u_1v_2{-}u_2v_1\bigr)\trans}
$u\times v\perp u,v$;\ $\|u\times v\|=\|u\|\|v\|\sin\theta$ (Fl\"ache).\\
\textbf{Spat}: $[u,v,w]=\inner{u}{v\times w}=\det[u\,v\,w]$ (Vol.).\\
$u\times v=-(v\times u)$.

\subsection{Geraden / Ebenen}
\textbf{Gerade}: $g:\,x=p+tr$.\\
\textbf{Ebene} (Pkt-Richt.): $x=p+su+tv$.\\
\textbf{Normalenform}: $\inner{n}{x-p}=0\Leftrightarrow n\trans x=d$.\\
\textbf{Koord.-Form}: $ax_1+bx_2+cx_3=d$.

\subsection{Abst\"ande}
\begin{itemize}
\item Pkt\,$q$--Ebene: $|n\trans q-d|/\|n\|$
\item Pkt\,$q$--Ger.\ ($\R^3$): $\|(q-p)\times r\|/\|r\|$
\item Windschief: $|\inner{p_2-p_1}{r_1\times r_2}|/\|r_1\times r_2\|$
\end{itemize}

\subsection{Projektion}
Spalten von $A$ lin.\ unabh.:
\fitline{\hat b=A(A\trans A)^{-1}A\trans b=Pb}
$P^2=P$,\ $P\trans=P$. Mit $Q$ orthonorm.: $P=QQ\trans$.\ $I-P$ projiziert auf Komplement.

\subsection{Quadriken}
$x\trans Ax+b\trans x+c=0$, $A$ symm.\\
\textbf{Klassifikation via EW-Signatur}:
\begin{tabularx}{\linewidth}{@{}l@{\ }X@{}}
alle $+$ & Ellipse / Ellipsoid \\
alle $-$ & leer \\
gemischt & Hyperbel / ein-\ \& zweischal.\ Hyp. \\
ein EW $=0$ & Parabel / Paraboloid \\
\end{tabularx}
\textbf{Hauptachsentrf.}: (1) lin.\ Term durch Verschiebung weg, (2) EVD $A=Q\Lambda Q\trans$, (3) $x=Qy\Rightarrow\sum\lambda_i y_i^2=c$.

% =================================================
\section{Vektorr\"aume}
% =================================================
\begin{itemize}
\item UVR: $0\!\in\!U$ \& abgeschl.\ unter $+$ und $\alpha\cdot$.
\item Basis = lin.\ unabh.\ + erzeugend; $\dim V$ = $|$Basis$|$ (eindeutig).
\item Steinitzscher Austauschsatz.
\item $f$ linear: $f(\alpha u+\beta v)=\alpha f(u)+\beta f(v)$.
\item Dimensionssatz: $\dim V=\dim\ker f+\dim\im f$.
\item $f$ inj.\ $\Leftrightarrow\ker f=\{0\}$;\ surj.\ $\Leftrightarrow\im f=W$.
\item Basiswechsel: $A'=T^{-1}AT$ ($\sim$ \"Ahnlichkeit).
\item Orthogonal: $Q\trans Q=I$, $|\lambda|=1$, $\det Q=\pm 1$.
\item 2D-Rotation: $\bigl(\begin{smallmatrix}\cos\theta&-\sin\theta\\ \sin\theta&\cos\theta\end{smallmatrix}\bigr)$.
\end{itemize}

% =================================================
\section{Klausur-Stolperfallen}
% =================================================
\begin{itemize}
\item LU: Multiplikatoren in $L$ nicht vergessen.
\item Cholesky: $\sqrt{\text{negativ}}\Rightarrow$ PD-Verstoss.
\item QR: vergewissern dass $Q$ orthonormal ($Q\trans Q=I$).
\item EVD: falsche alg.\ vs.\ geom.\ Vielfachheiten.
\item SVD: $\sigma_i$ niemals negativ; absteigend sortieren.
\item Char.\ Pol.\ Vorzeichen: $\det(A-\lambda I)$, nicht $(\lambda I-A)$.
\item $P^{-1}$ bei Diagonalisierung selbst berechnen, nicht raten.
\item Quadrik-Typ aus EW-Signatur, nicht aus $\det A$ allein.
\end{itemize}

\end{multicols*}
\end{document}
